\DocumentMetadata{
  pdfversion=2.0,
  pdfstandard=ua-2,
  lang=en-US,
  tagging=on,
  tagging-setup={math/setup=mathml-SE,tabsorder=structure}
}
\documentclass[11pt,letterpaper]{article}
\usepackage[margin=0.68in,includefoot]{geometry}
\usepackage{newcomputermodern}
\usepackage{amsmath,amscd,mathtools}
\providecommand{\square}{\mdlgwhtsquare}
\usepackage[hidelinks]{hyperref}
\hypersetup{
  pdftitle={Combinatorics PhD Exam, May 2023},
  pdfauthor={University of Florida Department of Mathematics},
  pdfsubject={Graduate mathematics examination},
  pdfdisplaydoctitle=true,
  bookmarksopen=true
}
\setcounter{secnumdepth}{0}
\setlength{\parindent}{0pt}
\setlength{\parskip}{0.28em}
\setlength{\emergencystretch}{3em}
\setlength{\leftmargini}{2.15em}
\setlength{\leftmarginii}{2.65em}
\setlength{\leftmarginiii}{3em}
\renewcommand{\labelenumi}{\textbf{\theenumi.}}
\pagestyle{empty}
\raggedbottom
\allowdisplaybreaks
\newenvironment{examproblems}[1]{%
  \begin{enumerate}%
  \setcounter{enumi}{#1}%
  \setlength{\topsep}{0.35em}%
  \setlength{\partopsep}{0pt}%
  \setlength{\itemsep}{0.48em}%
  \setlength{\parsep}{0.18em}%
}{\end{enumerate}}
\newenvironment{examsubparts}{%
  \begin{enumerate}%
  \setlength{\topsep}{0.18em}%
  \setlength{\partopsep}{0pt}%
  \setlength{\itemsep}{0.16em}%
  \setlength{\parsep}{0.1em}%
}{\end{enumerate}}
\newenvironment{examinstructions}{%
  \par\begingroup\itshape
}{%
  \par\endgroup\vspace{0.18em}
}
\makeatletter
\renewcommand\section{\@startsection{section}{1}{0pt}%
  {-1.25ex plus -.25ex minus -.1ex}%
  {0.55ex plus .1ex}%
  {\normalfont\large\bfseries}}
\renewcommand{\@maketitle}{%
  \newpage\null\vskip -1.2em%
  {\footnotesize\bfseries
    \href{https://www.ufl.edu/}{University of Florida}, \href{https://math.ufl.edu/}{Department of Mathematics}\par}%
  \vskip 0.18em%
  \tagstructbegin{tag=Title}%
    {\LARGE\bfseries\href{https://gma.math.ufl.edu/exams/combinatorics-phd/}{Combinatorics PhD Exam}, May 2023\par}%
  \tagstructend%
  \vskip 0.35em%
  \tagmcbegin{artifact}%
    \hrule height 1pt%
  \tagmcend%
  \vskip 0.55em%
}
\makeatother
\title{Combinatorics PhD Exam, May 2023}
\author{}
\date{}
\begin{document}
\maketitle
\thispagestyle{empty}
\begin{examproblems}{0}
\item
Let \(a_1,a_2,\ldots,a_{10}\) be positive integers not exceeding \(100\). Prove that there are disjoint, nonempty subsets~\(S\) and~\(T\) of \([10]\) so that 
\[
\sum_{i\in S}a_i=\sum_{j\in T}a_j.
\]
\item
Prove that the number of permutations of length \(n+1\) that have exactly two cycles is equal to the number of all cycles in all \(n!\) permutations of length~\(n\). (For instance, if \(n=3\), then both of these numbers are equal to \(11\).)
\item
Let~\(P\) be a convex polyhedron whose faces are all either~\(a\)-gons or~\(b\)-gons, and whose vertices are each incident to three edges. Let \(p_a\), \(p_b\), and~\(n\) respectively denote the number of~\(a\)-gonal faces, \(b\)-gonal faces, and vertices of~\(P\). Prove that \(p_a(6-a)+p_b(6-b)=12\).
\item
A city has three high schools, each of them attended by~\(n\) students. Each student knows exactly \(n+1\) students who attend a high school different from his. Prove that we can choose three students, one from each school, so that each of them knows the other two.
\item
Let \(L_n\) be the total number of leaves in all binary plane trees with~\(n\) vertices. These are trees in which all non-leaf vertices have one or two children, and every child is a left child or a right child, even if it is an only child. So \(L_1=1\), \(L_2=2\), and \(L_3=6\). Find an explicit formula for the numbers \(L_n\).
\item
Let \(c:S\to\{0,1\}^*\) be a prefix-free code in which \(b_i\) codewords have length~\(i\). Prove that \(\sum_i \frac{b_i}{2^i}\leq 1\).
\item
Let~\(T\) be a rooted non-plane~\(2\)-tree (that is, every non-leaf vertex has two children) with~\(n\) leaves. For instance, for \(n=4\), there are two such trees, one in which each leaf is at edge-distance two from the root, and one in which the edge-distances of the leaves from the root are~\(1\), \(2\), \(3\), and~\(3\). Let \(\operatorname{sym}(T)\) be the number of non-leaf vertices~\(v\) of~\(T\) such that the two children of~\(v\) are the roots of identical subtrees. Your answers to parts (a) and (b) can contain \(\operatorname{sym}(T)\), but your answer to part (c) should not.
\begin{examsubparts}
\item[\textnormal{(a)}]
How many automorphisms does~\(T\) have?
\item[\textnormal{(b)}]
How many different ways are there to bijectively label the leaves of~\(T\) with the numbers \(1,2,\ldots,n\)?
\item[\textnormal{(c)}]
Find an explicit formula for \(\sum_T \frac{1}{|\operatorname{Aut}(T)|}\), where the sum is taken over all rooted non-plane~\(2\)-trees with~\(n\) leaves. For example, for \(n=4\), the two trees mentioned above have, respectively, eight and two automorphisms, so the requested sum is \(\frac18+\frac12=\frac58\).
\end{examsubparts}
\item
Let~\(k\) be a fixed positive integer, and let \(F_k(z)=\sum_{n\geq k}c(n,k)z^n/n!\), where \(c(n,k)\) is the signless Stirling number of the first kind.
\begin{examsubparts}
\item[\textnormal{(a)}]
Find \(F_k(z)\) in an explicit form.
\item[\textnormal{(b)}]
Is \(F_k(z)\) an algebraic power series over~\(\mathbb{C}\)?
\item[\textnormal{(c)}]
Is \(F_k(z)\) a~\(d\)-finite power series over~\(\mathbb{C}\)?
\end{examsubparts}
\end{examproblems}
\end{document}
